\chapter{The OASIS Interface}

This interface was developped with \OA\ \citep{craig_2017} (see the module's
user guide:
\cite{valcke_oasis}).

\section{OASIS-related code}

Are listed below the sections of the code which are related to the use of OASIS:
\begin{itemize}
  \item the library \verb=oasislib=, needing only \verb=Share= and
        \verb=mpplib=;
  \item the module \verb=mod_oasis_interface.F90=, which links all
        \verb=oasislib= processes with the model's variables, and transfers the
        needed accessibility to \verb=oasislib= parameters (in the rest of the
        code, one only need to \verb=use mod_oasis_interface.F90=);
  \item several lines in \verb=mod_params.F90= for reading the \verb=oasisparam=
        namelist (and the parameter \verb=ioasiscpl= is added to the
        \verb=physicsparam= namelist);
  \item several lines in \verb=mod_regcm_interface.F90= for calling the needed
        subroutines along the run;
  \item several lines in \verb=regcm.F90= for implementing the MPI
        initialization of OASIS;
  \item appropriate lines in the \verb=Main='s \verb=Makefiles.am=;
  \item several lines in \verb=configure.ac= for the configuration options.
\end{itemize}

OASIS-related sections are enabled only if the builder is configured with the
option \verb=--enable-oasis=. In the code, these sections (except for the
parameter \verb=ioasiscpl= which is in \verb=physicsparam= anyways) are inserted
within
\begin{verbatim}
#ifdef OASIS
 ! Something to be done only in OASIS mode
#endif
\end{verbatim}

\section{Adding fields}

The interface is not exclusively for handling air-sea interactions, even though
it was
initially developped for it. Indeed, its types, variables, and subroutines
remain quite
general, such that one could for instance implement another exchange grid,
and/or set up
the fields which are needed for a precise coupling experiment.

The locations where lines should be added for implementing new fields are
notified with
the comment \verb=! OASIS field +++=, which may be \verb=grep=ped...

Any interaction between fields from OASIS and the model's variables should be
done in
\verb=mod_oasis_interface.F90=, where developpers are invited to use the model's
variables
which are already public.

Finally, simply adding a field should not need changes in the \verb=oasislib=,
although
consulting it may be useful for understanding then handling all types, variables
and
subroutines. Explicit comments are normally present in the code.

\section{eCLM land coupling}

With \verb=--enable-eclm= (which implies \verb=--enable-oasis=, and defines
\verb=ECLM=) the land surface model is eCLM, run as a separate OASIS
component:
\begin{itemize}
  \item BATS is not run. \verb=surface_model= in \verb=mod_lm_interface.F90=
        keeps only the tile values of the atmospheric forcing that BATS
        would set: the precipitation \verb=lms%prcp=, averaged over the
        surface step as in BATS, and the surface pressure \verb=lms%sfcp=.
        Fluxes, temperatures, albedos and roughness come from eCLM.
  \item The eCLM fields use their own grid, \verb=rlin=/\verb=rlim= in
        \verb=mod_oasis_interface.F90= (\verb=grdlnd=): the points of
        \verb=rcin=/\verb=rcim=, with the ocean masked instead of the land.
        Having separate names lets the ocean and the land couplings run
        together; the ocean grids keep their land mask.
  \item The received land fields are written on the land tiles of the
        \verb=lms= arrays with \verb=fill_land=; the wind speed at 10 m
        takes the direction of the wind at the lowest model level, and
        the drag is $|\tau|/|V|$, as used by the atmosphere.
\end{itemize}

\subsection{Multi-level fields}

Soil profiles, dust size bins and chemical species are exchanged as OASIS
bundles: \verb=type(info3dfld)= in \verb=mod_oasis_params.F90= holds the
number of members \verb=nz=, \verb=oasisxregcm_def_3dfield= declares the
variable with \verb=var_nodims = [1,nz]= on the 2d partition of its grid, and
\verb=oasisxregcm_rcv_3d= receives an array \verb=(j,i,nz)=. A bundle is one
\verb=namcouple= entry; OASIS interpolates each member in two dimensions
with the same weights, and both components must declare the same \verb=nz=.
A rank-3 call needs at least two members (OASIS3-MCT 5 aborts, version 6
takes it as an unbundled 3d field), so a single member is received as a
rank-2 array.

The sizes come from \verb=oasisparam=: \verb=oasis_nlevsoi= for the soil
profiles (\verb=RCM_TSOI=, \verb=RCM_H2OVOL=, \verb=RCM_H2OLIQ=,
\verb=RCM_H2OICE=), \verb=oasis_nddvel= for \verb=RCM_DDVEL=, the names in
\verb=oasis_megan_cmp= for \verb=RCM_FLXVOC=, and the parameter
\verb=oasis_ndust= (4) for \verb=RCM_FLXDST=. The soil profiles are stored
in \verb=lms%tsoi=, \verb=lms%sw_vol= and \verb=lms%h2osoi= (liquid plus
frozen water), allocated with \verb=oasis_nlevsoi= layers; the three
\verb=lms%sw= layers of the RegCM output and restart files keep the BATS
layout.

\subsection{Chemistry}

For the chemistry an ECLM build works as a CLM45 build: the code shared by
the two is under \verb=#if defined(CLM45) || defined(ECLM)=.
\begin{itemize}
  \item \verb=surface_model= averages over the subgrid the emissions
        \verb=lms%vocemiss=, \verb=lms%dustemiss= and the deposition
        velocities \verb=lms%ddepv= into \verb=voc_em_clm=,
        \verb=dustflx_clm= and \verb=ddepv_clm= of
        \verb=mod_atm_interface.F90=, which the chemistry sees as
        \verb=cvoc_em_clm=, \verb=cdustflx_clm= and \verb=cddepv_clm=.
  \item Isoprene: \verb=RCM_FLXVOC= is a bundle of MEGAN compounds; the one
        named \verb=ISOP= in \verb=oasis_megan_cmp= goes to the isoprene
        tracer. \verb=rcm_megan_enabled= is set when it is coupled, and the
        isoprene of the inventory is then replaced (\verb=mod_che_emission=).
  \item Methane: \verb=RCM_FLUXCH4= (carbon) is converted to methane mass
        and added to the anthropogenic emission.
  \item Dust: with \verb=ichdustemd = 3= the dust flux of eCLM
        (\verb=RCM_FLXDST=, 4 bins) is redistributed on the dust bins of
        RegCM (\verb=clm_dust_tend=). The RegCM dust scheme instead uses the
        fraction of ground emitting dust \verb=RCM_WT= (\verb=lms%wt=).
  \item The biogenic NO$_x$ emission (\verb=ichbion = 1=) takes the soil
        moisture of the second and the temperature of the first and third
        layers of \verb=RCM_H2OVOL= and \verb=RCM_TSOI=; the soil water of
        the top 10 cm, \verb=RCM_H2O10CM=, is \verb=ssw2da=.
  \item \verb=dustflx_clm=, \verb=voc_em_clm= and the soil profiles
        \verb=tsoi=, \verb=sw_vol= are in the restart file, the profiles on
        their own dimension \verb=eclm_soil_layers=: the hydrostatic and
        non-hydrostatic cores compute the chemistry before the surface.
  \item \verb=RCM_DDVEL= is received, but \verb=ddepv_clm= is not used by
        the dry deposition, as for CLM45.
\end{itemize}

\subsection{OASIS3-MCT versions}

The routines and arguments used by the interface are the same in
OASIS3-MCT 5.0, in the J\"ulich patched 5.0 used with eCLM in TSMP, and in
OASIS3-MCT 6. The J\"ulich patches rename the MCT
modules (to coexist with the MCT of CLM in one executable) and add
ensemble support for PDAF; neither changes the interface used by RegCM.
